\section{Cuántas componentes puede tocar un término}
Conservamos siempre el mismo $D_0=A\setminus S$. Para cada fibra RC distinguida $b$
la masa $y_b$ es suma de las $x_i$ incidentes, todos en una sola componente.
En un átomo RC compatible $y_b(X)$ es binaria, también cuando la fibra tiene varios sitios.
Las fórmulas del exterior con $f_d(y)=(1-y)^d-1$, $E_{ab}=\int z_x^2\,dM_a^R dM_b^C$, dan
\[
 R=\text{afín}+s^{-1}R_0(sy)+E_{22}/2+s\Psi_R,\qquad
 A_*=A_0(sy)+(E_{22}-E_{12}-E_{21})/2+s\Psi_A .
\]
El campo afín y las constantes se extraen con la misma convención en todas las caras.
Se da una masa independiente $h_g$ a las copias RC de cada componente en las medidas
$M_d=\mu^0+\sum_b h_{g(b)}f_d(y_b)\delta_b$. Se conserva el balance porque cada sitio
aporta una vez a cada eje en su componente. El núcleo $z_x$ depende de los $h_{g(b)}y_b$.
Las fórmulas holomorfas previas se expanden en los discos fijos para $v$ fijo.
Así a orden $s^l$, para $l\ge1$, sólo aparecen los siguientes tipos de términos.
\begin{enumerate}
\item Ordinarios: polinomios de grado $\le l+1$ en la acción y $\le l$ en el prefactor.
Incluyen $R_0,A_0$, las medidas sin marcas y con una marca en $E$. Con una sola se obtiene
$s f_2$ en la acción y $-s f_1$ en el prefactor, multiplicados por el núcleo desarrollado.
\item Dos masas en $E$: una sola amplitud común para el siguiente par de términos,
multiplicado por $s^2$ y por el Taylor del mismo $z_x^2/2$. Sea $F=f_2$, $L=f_1$.
En la acción aparece $P=F_iF_h$, y en el prefactor $B=F_iF_h-L_iF_h-F_iL_h$.
Los índices aquí son fibras marcadas de fila y columna. En orden $l$ el multiplicador
tiene grado $\le l-2$ en las masas lineales $y$. Por tanto su soporte tiene a lo sumo
$l$ grupos aun contando $i,h$. Esta pareja se marca siempre conjuntamente.
\item Un término de $s\Psi_R,s\Psi_A$ o $d_T$ toca a lo sumo $l-1$ grupos a orden $l$.
Para $\Psi$ una componente de masa cero deja de influir en sus integrales; cada componente
presente cuesta un paso $h_g=s$, además del $s$ inicial. Para $d_T$ se acaba de demostrar
lo mismo en los átomos. Los términos constantes comunes se pueden quitar.
\end{enumerate}

Marcamos cada interacción con un parámetro de Taylor para extraer las contribuciones
mixtas, polarizando las repeticiones. En el valor exacto y en el canónico un conjunto
de interacciones que se parte en dos conjuntos sobre grupos separados tiene coeficiente
mixto cero: la partición factoriza, la silla se descompone y la matriz del determinante
es diagonal por bloques. La misma propiedad vale para las derivadas en campos locales.
Por tanto el grafo de incidencia entre los grupos que intervienen y las interacciones
del monomio debe ser conexo. El soporte cuenta dependencias en las activaciones, no en los
coeficientes exteriores congelados, comunes a todos los casos. En orden cero, el error
basal es una suma de valores locales de un solo grupo.

Sean $h_a$ el número de grupos de una interacción, $l_a$ su coste y
$\eta_a=l_a+1-h_a\ge0$. Sea $L_I$ el número de ciclos independientes de aquel grafo
bipartito, con una incidencia por cada grupo distinto en cada interacción.
Para un coeficiente conexo de orden positivo $w$ su número $d$ de grupos satisface
\[
 d=1+\sum_a(h_a-1)-L_I=w+1-\sum_a\eta_a-L_I .
\]
Los posibles coeficientes mixtos se estudian en el origen de sus amplitudes con campos
locales libres; el teorema implícito y la expansión de las particiones finitas dan
restos uniformes a orden fijo en un polidisco menor. Esto no interpola la sucesión de
enteros $j$. Si un término tiene soporte menor que $c$, el alternar sobre todas las
activaciones lo mata sumando sobre una componente ausente.
